\printedpage{24}
% head: A. Grothendieck : Dérivateurs, Ch. VI, Édition des Universités de Montpellier II et Paris VII

\msp{page 41}

\begin{flushright}
15.11.90
\end{flushright}

I think I finally understood, last night, the exact situation. There is reason, in Cat$/S$ and in the other categories of the diagram p. 8 (all mapping to Cat$/S$) to consider four sorts of ``$W$-equivalences''. % [sic] qu j'ai -> que j'ai, letter missing

\begin{itemize}
\item[a)] The \textit{``coarse'' $W$-equivalence} : \begin{tikzcd}[baseline=(s.base), row sep=small, column sep=small]
f : X \arrow[rr] \arrow[dr] & & Y \arrow[dl] \\
& |[alias=s]| S &
\end{tikzcd} in Cat$/S$ is a coarse $W$-equivalence if as an arrow of Cat, it is a $W$-equivalence. (\footnote[17]{stable under base change both smooth and proper, more generally if $S' \longrightarrow S$ is a Hot-fibration.}).

\item[b)] The \textit{``fiberwise'' $W$-equivalence} (over $S$) : $\forall\ s \in S$, $f_s : X_s \longrightarrow Y_s$ is a $W$-equivalence. (\footnote[18]{stable under any base change $S' \longrightarrow S$}). % [sic] chagement -> changement, letter missing

\item[c)] The \textit{$W$-equivalence ``local over $S$''} : $\forall\ s \in S$, $f/s : X/s \longrightarrow Y/s$ is a $W$-equivalence. Means that it is a $W$-equivalence, and remains so under any base change $S' \longrightarrow S$ which is a local isomorphism. (\footnote[19]{should be stable under \textit{smooth} base change $S' \longrightarrow S$}).

\item[c$'$)] The \textit{$W$-equivalence ``colocal over $S$''} : $\forall\ s \in S$, $s\backslash f : s\backslash X \longrightarrow s\backslash Y$ is a $W$-equivalence. Means that it is a $W$-equivalence, and remains so under any base change $S' \longrightarrow S$ which is a local co-isomorphism. (\footnote[20]{should be stable under \textit{proper} base change over $S$}).
\end{itemize}

I am going to summarize in a recapitulative statement the relations between these four notions of $W$-equivalence, thinking of the case $W = W_\infty$. Once this is done, I shall concern myself with isolating the minimal conditions on a general $W$ which ensure that it is true for $W$.

\textbf{Proposition 1.} \textit{Let} \begin{tikzcd}[baseline=(s.base), row sep=small, column sep=small]
f : X \arrow[rr] \arrow[dr] & & Y \arrow[dl] \\
& |[alias=s]| S &
\end{tikzcd} \textit{be in} Cat$/S$.

\begin{itemize}
\item[1)] \textit{Local $W$-equivalence over $S$, and colocal $W$-equivalence over $S$, imply coarse $W$-equivalence :}
\[
\begin{tikzcd}
\text{c} \arrow[dr, Rightarrow] & & \text{c}' \arrow[dl, Rightarrow] \\
& \text{a} &
\end{tikzcd}
\]

\printedpage{25}
% head: A. Grothendieck : Dérivateurs, Ch. VI, Édition des Universités de Montpellier II et Paris VII

\item[2)] \textit{If $X$, $Y$ are proper over $S$, or $X$, $Y$ smooth over $S$, or $X$, $Y$} Hot-\textit{fibered over $S$, then fiberwise $W$-equivalence implies that it is a} universal \textit{$W$-equivalence over $S$ ($W$-equivalence, and remains so under} any \textit{base change $S' \longrightarrow S$), hence in these three cases, we have}
\[
\begin{tikzcd}
& \text{b} \arrow[dl, Rightarrow] \arrow[dr, Rightarrow] & \\
\text{c} \arrow[dr, Rightarrow] & & \text{c}' \arrow[dl, Rightarrow] \\
& \text{a} &
\end{tikzcd}
\]
(\footnote[21]{It suffices $X$, $Y$ $W$-proper over $S$, or $W$-smooth over $S$, or $W$-fibered over $S$.}).

But we can be more precise about each of these three cases :

\item[(2 fib)] \textit{If $X$, $Y$ are} Hot-\textit{fibered over $S$, then coarse $W$-equivalence of $f$ implies that it is a fiberwise $W$-equivalence, hence the four notions of $W$-equivalence are equivalent for $f$.}
\[
\text{a} \Longleftrightarrow \text{c} \Longleftrightarrow \text{c}' \Longleftrightarrow \text{b}
\]
(\footnote[22]{a $\Longrightarrow$ b only assuming L5, see below p. 43.})

\item[(2 smooth)] \msp{page 42}\textit{If $X$, $Y$ are smooth over $S$, then fiberwise $W$-equivalence over $S$ is equivalent to colocal $W$-equivalence over $S$ and implies} universal \textit{$W$-equivalence over $S$, hence in this case} % [sic] lisse -> lisses, plural agreement with X, Y
\[
\begin{tikzcd}
\text{b} \arrow[r, Leftrightarrow] \arrow[d, Rightarrow] & \text{c}' \\
\text{c} \arrow[d, Rightarrow] & \\
\text{a} &
\end{tikzcd}
\]

\printedpage{26}
% head: A. Grothendieck : Dérivateurs, Ch. VI, Édition des Universités de Montpellier II et Paris VII
(\footnote[23]{it even suffices $X$, $Y$ $W$-smooth over $S$.}).

\item[(2 prop)] \textit{If $X$, $Y$ are proper over $S$, then fiberwise $W$-equivalence over $S$ is equivalent to local $W$-equivalence over $S$ :}
\[
\begin{tikzcd}
\text{b} \arrow[r, Leftrightarrow] \arrow[d, Rightarrow] & \text{c} \\
\text{c}' \arrow[d, Rightarrow] & \\
\text{a} &
\end{tikzcd}
\]
(\footnote[24]{It even suffices $X$, $Y$ $W$-proper over $S$.}).
\end{itemize}

\textsc{Proof} 1) is none other than L4 (\footnote[25]{No, there is confusion, L4 is weaker, cf. p. 44 the axiom L4 bis.}) and its dual L4$'$ for a basic localizer $W$, and we have seen that if we assume L1, L3, then L4 is equivalent to L4$'$, and in the notion of basic localizer we shall always assume that this is satisfied. Also, we assume that L3 is true even in the stronger form L3 quater (p. 1) : $W_\infty \subseteq W$, i.e. every Hot-equivalence is a $W$-equivalence.

(2 smooth) If $X$, $Y$ \texttt{[are]} smooth over $S$, then the $X_s \longrightarrow s\backslash X$, $Y_s \longrightarrow s\backslash Y$ are $W_\infty$-equivalences, hence $W$-equivalences (it would even suffice that $X$, $Y$ be $W$-smooth$/S$), hence the commutative diagram
\[
\begin{tikzcd}
X_s \arrow[r, "f_s"] \arrow[d, "i_s^X"'] & Y_s \arrow[d, "i_s^Y"] \\
s\backslash X \arrow[r, "s\backslash f"] & s\backslash Y
\end{tikzcd}
\]
shows us that $f_s$ is a $W$-equivalence if and only if $s\backslash f$ is, whence b $\Longleftrightarrow$ c$'$. To see that fiberwise $W$-equivalence is even a \textit{universal} $W$-equivalence over $S$, we are reduced to showing that it is a coarse $W$-equivalence (since smoothness \texttt{[is]} stable under base change over $S$), and we win since c$'$ $\Longrightarrow$ a by 1).

The case (2 proper) is dual.

(2 fib) If $X$, $Y$ are Hot-fibered over $S$, then
\printedpage{27}
% head: A. Grothendieck : Dérivateurs, Ch. VI, Édition des Universités de Montpellier II et Paris VII
\msp{page 43}
the inclusions
\[
X_s \longrightarrow s\backslash X, \qquad X_s \longrightarrow X/s,
\]
and likewise for $Y$, are $W_\infty$-equivalences, hence $W$-equivalences (it would even suffice that $X$, $Y$ be $W$-fibrations over $S$ (\footnote[26]{compare def. 3, p. 117, in (I).}), i.e. by definition that the preceding functors be $W$-equivalences, and that the same condition be satisfied after any base change over $S$), the argument in the smooth resp. proper case already shows that
\[
\begin{tikzcd}
\text{c} & \text{b} \arrow[l, Rightarrow] \arrow[r, Leftrightarrow] \arrow[d, Rightarrow] & \text{c}' \\
& \text{a,} &
\end{tikzcd}
\]
and that fiberwise $W$-equivalence (or local $W$-equivalence over $S$, resp. colocal over $S$) already implies \textit{universal} $W$-equivalence over $S$. (\footnote[27]{valid assuming only $X$, $Y$ $W$-fibered over $S$.})

The question a $\Longrightarrow$ b is on the other hand distinctly more delicate, and we have seen that this is by no means the case for every basic localizer - probably not for the one associated to the derivator Der($\mathbf{Z}$-Mod), in any case not for the one associated to Der($k$-Mod), when 2 is invertible in $k$. (It would even suffice that $\exists\ n \in \mathbf{N}$, $n \geq 2$, such that $n$ invertible in $k$ - the typical counterexample is $X = S^1 \xrightarrow{\ n\ } S^1 = Y = S$, the étale covering of order $n$ of the circle.)

The implication a $\Longrightarrow$ b is thus a very strong condition on $W$, I am going to denote it L5 (the old L5, cf. page 2, not deserving to survive, for it is a botched compromise).

\textbf{L5} \textit{Let} \begin{tikzcd}[baseline=(s.base), row sep=small, column sep=small]
X \arrow[rr, "f"] \arrow[dr] & & Y \arrow[dl] \\
& |[alias=s]| S &
\end{tikzcd} \textit{be in} Cat\textit{, with $X$, $Y$} Hot-\textit{fibered over $S$. If $f$ is a $W$-equivalence (``coarse'' relative to $S$), then it is a fiberwise $W$-equivalence} (\footnote[28]{In XII (p. 12) we see that if $W$ satisfies L1, L3 quater, L4 bis (cf. next page), L5, then \fbox{$W = W_\infty$}. So it is perfectly illusory, if we assume L5 to go on dragging a $W$ along.}).

We have seen that up to an $\varepsilon$ (``lemma'' p. 23), L5 is a consequence of

\textbf{L6} \textit{Let $f : X \longrightarrow Y$ be a $W$-equivalence, $Y' \longrightarrow Y$ a perfect morphism (i.e. smooth and proper), then $f' : X' = X \times_Y Y' \longrightarrow Y'$ is a $W$-equivalence.}

\msp{page 44}
Let us note the quite natural variant

\textbf{L6 bis} \textit{As in L6, but assuming only $Y' \longrightarrow Y$ a} Hot-\textit{fibration.}

\printedpage{28}
% head: A. Grothendieck : Dérivateurs, Ch. VI, Édition des Universités de Montpellier II et Paris VII

There is a chance that in practice, when L6 (or only L5) is satisfied, L6 bis is too. % [sic] verifié -> vérifié, missing accent

We are going to introduce an axiom

\textbf{L7} \ a) \textit{Let $(f_i : X_i \longrightarrow Y_i)_{i \in I}$ be a family of arrows in} Cat\textit{, let $f : \coprod X_i \longrightarrow \coprod Y_i$ be the corresponding arrow on the sum categories. For $f \in W$, it is necessary and sufficient that $f_i \in W$ for every $i \in I$.}

\hspace*{2em}b) \textit{Let $f : X \longrightarrow Y$ be in $W$. If $Y$ is} $0$\textit{-connected, $X$ too.}

This axiom is equivalent to the conjunction of La ($f \in W \Longrightarrow \pi_0 f$ bijective) and the validity of the ``Lemma'' page 23, which I put here again as a reminder :

\textbf{Lemma.} \textit{The preceding axiom L7 on the basic localizer $W$ is equivalent to the following condition : $f : X \longrightarrow Y$ is in $W$ if and only if $\pi_0(f)$ is bijective, and $f$ induces $W$-equivalences $f_i : X_i \longrightarrow Y_i$ on the connected components.}

I must come back once more to 1) of the proposition, it is not true that ``it is'' the condition L4 on $W$ (which I assume, I said), it is L4 in the particular case where $Y = S$. On reflection, I see it as a strengthened variant of L4, which does not seem likely to follow from L4 :

\textbf{L4 bis} \textit{Let} \begin{tikzcd}[baseline=(s.base), row sep=small, column sep=small]
X \arrow[rr, "f"] \arrow[dr] & & Y \arrow[dl] \\
& |[alias=s]| S &
\end{tikzcd} \textit{be in} Cat\textit{, let us suppose that $f$ is a local $W$-equivalence over $S$ (i.e. $f/s : X/s \longrightarrow Y/s$ a $W$-equivalence $\forall\ s \in S$), then $f$ is a $W$-equivalence} (\footnote[29]{equivalent to the dual axiom, in terms of the $s\backslash f$}).

\msp{page 45}
In what follows, by a ``basic localizer'', we shall mean a set
\[
W \subseteq \mathrm{Fl}\,\mathrm{Cat}
\]
satisfying (at least) the axioms % [sic] axioms -> axiomes, French spelling
\[
\text{L1, L3 quater, L4 bis, L7,}
\]
i.e.

\begin{itemize}
\item[a)] \textit{Saturation} of $W$ (conditions a) b) c) p. 1)

\item[b)] $W_\infty \subseteq W$ $[\subseteq W_0]$ (\footnote[30]{I am no longer sure whether it is really useful to impose $W_\infty \subseteq W$. Take instead the final object axiom, namely b$'$). If we want a) b$'$) c) + L5, then we shall have $W \subseteq W_\infty$. \textbf{NB} a) b) c) + L5 implies $W = W_\infty$, cf. annotation page 43.}) (\footnote[31]{cf. d)}).

\item[c)] If $f : X \longrightarrow Y$ over $S$ is locally over $S$ a $W$-equivalence, then $f \in W$ (L4 bis)
\end{itemize}